% Part II, discussion (included by main_en.tex).
% DRAFT, written with 8 of the 12 GRPO runs; to be checked against the final numbers.

\section{Discussion}\label{sec:discussion}

\paragraph{What sharing the backbone buys.} The judge on the shared backbone gives the same answers as the judge
loaded on its own, and the trainer holds one copy of the backbone instead of two: \NUM{mem_saving_load} less GPU memory
after loading. For a 4B policy on a 96\,GB GPU this is a convenience; on a smaller GPU, or with a larger model, it can
decide whether a judge fits next to the policy at all. It costs \NUM{adapter_switch_ms}\,ms per adapter switch and
some speed from applying the judge's adapter unmerged.

\paragraph{Why the judge adds nothing on top of the verifier.} The verifier gives every completion an exact verdict,
and with $\lambda<1$ the judge can only reorder completions with the same verdict. That leaves two places for it to
help: among wrong completions, by preferring the ones closer to right, and in groups where every completion is wrong,
where the verifier gives no signal at all. After the first steps few groups are all wrong
(\NUM{allwrong_rlvr} over the last 20 steps of \texttt{rlvr}), and among same-verdict completions a judge's preference
can be about style as much as about correctness: under the gate the policy wrote longer solutions that repeat
themselves, without becoming more accurate (\S\ref{sec:res-repeat} tests whether the judge's score rises with
repetition). % TODO: rewrite with the probe's result. Our runs do not rule out a
gain with another weight $\lambda$, harder problems with more all-wrong groups, a stronger judge, or longer training;
they show that at this scale the gain, if any, is below the resolution of 2--3 seeds and the evaluation noise.

\paragraph{A judge is not a verifier.} Offline, \kev{}'s $\pc$ separates correct from incorrect base-policy solutions
with an AUROC near 0.9 (\S\ref{sec:res-study}). The probes showed that it can be moved by edits that do not change the
answer, if only slightly (a claim that every step was verified). Trained against $\pc$ alone, the policy found
solutions the judge overrates: its wrong answers reach a $\pc$ the judge used to give only to right ones, the judge's ranking within a group decays towards chance, and the
accuracy on every benchmark but GSM8K falls below the base model's. This is the reward over-optimisation that
\citet{gao2023scaling} describe, at small scale and within 200 steps. An offline AUROC measures a judge on the
distribution it was trained for, not on the distribution a policy optimised against it will produce.

\paragraph{Limitations.} One policy and judge size (4B), one task (MATH, evaluated on MATH-500, GSM8K and three
competition sets), 200 steps, one learning rate chosen for \texttt{rlvr}, one judge weight, and 2--3 seeds per
condition. The intervals resample problems and not seeds. Greedy evaluation on vLLM varies by about
\NUM{szsd_math500} points between repeated evaluations of the same model (\S\ref{sec:res-cond}). Other experiments shared
the GPUs during training, so we report no training times except those measured on an idle GPU.

\paragraph{Relation to the earlier versions.} Versions 1--3 of this paper reported a leaderboard on which \kev{}'s
reward improved on the verifier, and memory savings computed from an assumed factor. Neither was measured (see the
notice at the top and \file{CORRECTION.md}). The measured answer: the shared backbone saves one copy of the backbone
and leaves the judge unchanged, and adding \kev{}'s judgment to the verifier did not improve GRPO on MATH in our runs.
